% Structure follows regular ReScience C replication articles: Introduction, Scope of
% reproducibility, Methodology, Results, Discussion, Conclusion.
%
% Every number here comes from a report script run over a committed result file
% (`make figures` prints them all), never typed by hand.
%
% Sources: raw observations and the conditions stated before each run are in
% OVERNIGHT_LOG.md; departures from the papers are in DEVIATIONS.md.

% Figures are not committed. `make figures` rebuilds them in about two minutes without a
% GPU. This guard keeps the paper compilable before that has been run.
\newcommand{\resfig}[3]{%
  \begin{figure}[h]\centering
  \IfFileExists{#1}{\includegraphics[width=\linewidth]{#1}}%
    {\fbox{\parbox{0.9\linewidth}{\centering\vspace{1em}\small\texttt{\detokenize{#1}}\\
     not built. Run \texttt{make figures} in the repository root.\vspace{1em}}}}
  \caption{#2}\label{#3}\end{figure}}

\section{Introduction}

A Neural ODE replaces the discrete layers of a network with a learned differential
equation\supercite{chen2018neural}. The hidden state $h(t)$ evolves according to
$\dot h = f_\theta(h,t)$, and the output is the state at $t=1$, computed by a numerical
solver. Depth becomes integration time. Memory can be made constant in depth through the
adjoint method. The number of function evaluations, written NFE, measures how hard the
learned dynamics are to integrate.

Dupont, Doucet and Teh\supercite{dupont2019augmented} show that this construction has a
structural limitation. The flow of an ODE is a homeomorphism, so trajectories cannot cross
and the map from input to terminal state cannot change the topology of the data. A Neural
ODE therefore cannot represent functions that require such a change. When trained on one, it
does not fail gracefully. It makes its flow steeper, which raises the NFE. The proposed
remedy is augmentation: append zero-valued dimensions to the state before integrating, so
that the flow can separate the data in a higher-dimensional space. The paper supports this
with a proposition, experiments on nested geometries, and image benchmarks at matched
parameter counts.

These claims are worth replicating for two reasons. They are structural rather than
statistical, so they should hold or fail sharply. They are also cheap to test, which puts
them within reach of a single GPU. The central empirical signature, the NFE, is a property of
the solver as much as of the model. An adaptive solver given too loose a tolerance returns a
plausible trajectory and a plausible NFE without tracking the flow. The vector field of a
Neural ODE also stiffens during training, so a tolerance that is adequate early need not stay
adequate. Several conclusions reached during this work were withdrawn for this reason
(Section~\ref{sec:withdrawn}). We
therefore attach a reconstruction check to every measurement and quote each quantity at a
tolerance where that check passes. Section~\ref{sec:recon} defines the check.

We reimplemented both models from the descriptions in the papers, without consulting either
author's code, and tested eight claims. Four are from Dupont et al.\ and four are the solver
and memory claims of Chen et al.\ on which they rest. The central claims reproduce. The
one-dimensional impossibility result reproduces in a strong form: the Neural ODE fails at the
error floor the proposition predicts. The cost and generalisation signatures reproduce on two
geometries, and we extend the generalisation experiment into a systematic sweep. On the image
benchmarks the ordering reproduces on both datasets, while two of the four absolute
accuracies fall below the published values. Of the four supporting claims, two reproduce, one
reproduces partially, because a condition stated in advance is violated at tolerances where the
solver is not integrating, and one does not reproduce. We do not attempt the benchmark table of Chen
et al.

Two earlier reports examined the same paper for the NeurIPS 2019 Reproducibility Challenge.
\citet{lehnert2019reproducibility} reproduced most of its results using the authors' released
implementation, and added ablations with fixed-step solvers. \citet{ho2019reproducibility}
reimplemented the image models and recovered the ordering of the two, at accuracies about one
point below the original on MNIST and about eight below on CIFAR-10, without accounting for the
difference. Both reports were rejected by the challenge and remain unpublished. This replication
asks different questions. The first is whether the published description alone suffices, without
reference to the authors' code. The second is whether the reported measurements survive a check
that the solver integrated the flow. The third is whether the four claims of Chen et al.\ on
which the argument rests hold as well. Where our own accuracies fall short, we say by how much
and what we think explains it.

\section{Scope of reproducibility}
\label{sec:scope}

We test the following claims. Claims 1 to 4 are from Dupont et al. Claims 5 to 8 are from
Chen et al.

\begin{enumerate}
  \item \textbf{A one-dimensional Neural ODE cannot represent the crossing map.} The flow of
        a scalar ODE preserves order, so no Neural ODE can approximate $g(x)=+1$ for $x<0$
        and $g(x)=-1$ for $x>0$. One augmented dimension is sufficient. (Proposition 1,
        Figure 3.)
  \item \textbf{The cost of a Neural ODE grows with the training budget, and an augmented
        model stays flat.} On a task that requires the flow to separate nested classes, the
        NFE of the Neural ODE increases with the budget. The augmented model is cheaper and
        flatter. (Section 4.2, Figure 6.)
  \item \textbf{Neural ODEs generalise poorly across a region absent from training.}
        Removing an angular wedge from the training set produces a large gap on that wedge
        for the Neural ODE. The augmented model generalises across it. (Section 5.1,
        Figure 9.)
  \item \textbf{At matched parameter counts, augmentation improves image accuracy and lowers
        cost.} On MNIST and CIFAR-10 the augmented model is more accurate than a
        parameter-matched Neural ODE, at lower NFE. (Table 1.)
  \item \textbf{Solver error falls and cost rises as the tolerance tightens}, and forward
        computation time is approximately proportional to NFE. (Chen et al., Figure 3a--b.)
  \item \textbf{The backward pass costs about half the forward pass in function
        evaluations.} (Chen et al., Figure 3c.)
  \item \textbf{The number of function evaluations increases during training.} (Chen et al.,
        Figure 3d.)
  \item \textbf{The memory cost of the adjoint method is constant in the effective depth of
        the model.} (Chen et al., Table 1, memory column.)
\end{enumerate}

Claims 1 to 3 are qualitative. Claims 4 to 8 are quantitative. For the quantitative claims
the tolerance at which the ODE is solved determines whether a measurement is meaningful,
which is the subject of Section~\ref{sec:recon}.

We do not attempt the benchmark table of Chen et al. Our MNIST models are our own seeded
baselines and are not comparable to that table in architecture, budget or solver. We also
exclude the Latent ODE and ODE-RNN models of Rubanova et
al.\supercite{rubanova2019latent}, and the SVHN and ImageNet experiments of the original
paper. None of these is needed to test the claims above, and all exceeded the compute
available to us.

\begin{table}[h]
  \centering
  \small
  \caption{The eight claims, their location in the original articles, and the outcome.}
  \label{tab:claims}
  \begin{tabular}{@{}clll@{}}
    \toprule
    \# & Claim & Original & Outcome \\
    \midrule
    1 & 1-D crossing map unrepresentable & Dupont, Prop.~1, Fig.~3 & reproduced \\
    2 & NFE grows with training budget   & Dupont, Fig.~6         & reproduced \\
    3 & Gap across an unobserved region  & Dupont, Fig.~9         & reproduced \\
    4 & Matched-parameter image accuracy & Dupont, Table~1        & partial \\
    5 & Solver error and cost vs tolerance & Chen, Fig.~3a--b     & partial \\
    6 & Backward NFE about half forward  & Chen, Fig.~3c          & not reproduced \\
    7 & NFE grows during training        & Chen, Fig.~3d          & reproduced \\
    8 & Adjoint memory constant in depth & Chen, Table~1          & reproduced \\
    \bottomrule
  \end{tabular}
\end{table}

\section{Methodology}
\label{sec:methods}

\subsection{Model descriptions}

All models integrate a learned vector field $f_\theta$ from $t=0$ to $t=1$ and classify the
terminal state with a single linear map. The linear readout is required by the claims under
test. A nonlinear head could separate the classes by itself, which would make the properties
of the flow irrelevant. The one exception is claim 1, whose model has no readout at all. Its
output is the data coordinate of the terminal state. A learnable readout $w\,\phi(x)+b$ with
$w<0$ would reverse the order that the flow preserves, so a one-dimensional model could
represent the crossing map through the readout alone. Proposition 1 is a statement about the
flow, and its error floor applies to the flow without a readout. An augmented model, written ANODE, appends $p$ zero-valued dimensions
to the state before integration. Setting $p=0$ gives the plain Neural ODE, written NODE.

For the one- and two-dimensional tasks, $f_\theta$ is a multilayer perceptron of width 32
that takes the state and the time, following Appendix F.1.1 of the original paper. Claim 1
also uses the training settings Appendix F.2.1 gives for $d=1$: 3000 points, batch size 64,
learning rate $10^{-3}$ and 50 epochs, with $p=1$ and with $p=5$, the value the original found
best. It is applied
directly to the raw input. No encoder precedes the flow, because a learned encoder is itself
a diffeomorphism and would supply the warp that the topological argument assumes is absent.

For the image tasks, $f_\theta$ is the convolutional field of Appendix F.1.2, a stack of
$1\times1$, $3\times3$ and $1\times1$ convolutions. The time is injected as an additional
channel before every convolution. Augmentation adds $p$ zero channels to the input. Capacity
is matched between NODE and ANODE through the number of filters, as in the original. On MNIST
we use 92 filters (85{,}316 parameters) against 64 filters with $p=5$ (85{,}462). On CIFAR-10
we use 125 filters (173{,}611) against 64 filters with $p=10$ (172{,}452). These counts are 0.4\% to 1.1\% above the counts stated in the original Table 1. We attribute the difference to
the flattening readout and record it in Section~\ref{sec:deviations}. The harness asserts that
the two arms are matched to within 2\% before every run.

Both models were written from the descriptions in the two papers. We did not consult either
author's implementation. The original code for Figure~3 of Chen et al., which its authors
later shared, was read only for the diagnosis of claim 6 (Section~\ref{sec:communication}). The solver library \texttt{torchdiffeq}\supercite{torchdiffeq} is
used only to integrate. Its adaptive Dormand--Prince method, \texttt{dopri5}, is the default
throughout, except where a claim concerns the solver itself.

\subsection{Datasets}

The toy task is the concentric-spheres dataset of Appendix F.2.1. A filled inner disk
($\lVert x\rVert \le 0.5$) is enclosed by an annulus ($1.0 \le \lVert x\rVert \le 1.5$),
sampled in a $1{:}2$ ratio. The inner class is surrounded by the outer one, so no flow in the
plane can separate them without trajectories crossing. The geometry is what makes the task
hard. Where stated, we repeat an experiment on a thinner concentric-circles geometry as a
robustness check.

The image tasks use MNIST and CIFAR-10, scaled to $[0,1]$, with no data augmentation. This
keeps the comparison about the flow rather than about an augmentation pipeline. Neither
dataset is redistributed with our code. The loaders fetch them on first use.

\subsection{Solver tolerance and the reconstruction check}
\label{sec:recon}

An adaptive ODE solver takes a tolerance and returns a trajectory and the number of function
evaluations it used. Neither output reports failure. If the tolerance is too loose for the
field, the solver still returns a state and an NFE, and any quantity derived from them
describes the solver rather than the model. The vector field of a Neural ODE stiffens during
training, so a tolerance fixed at the start of a run can stop being adequate before the end.

We therefore treat the tolerance as an experimental axis and attach a check to every
measurement. The flow of an ODE is invertible, so integrating forward from $0$ to $T$ and
back to $0$ must return the input. We integrate in both directions with the same solver and
tolerance and record the relative reconstruction error
\[
  \varepsilon_{\text{recon}}
  \;=\;
  \frac{\lVert \phi_{T\rightarrow 0}(\phi_{0\rightarrow T}(x)) - x \rVert}{\lVert x \rVert}.
\]
A measurement passes the check when $\varepsilon_{\text{recon}} < 10^{-2}$. We call such a
tolerance a checked tolerance. The check is a necessary condition rather than a sufficient
one. It screens for a gross failure to track
the flow, and it is cheap enough to run on every measurement. A flow that does not return to its
starting point has not been integrated. One that does return may still be inaccurate.

Two reporting rules follow. First, a quantity compared between two models is quoted at the loosest tolerance at which both
models pass the check on every seed. This prevents either model from being favoured by a
tolerance that suits only it. Second, measurements that fail the check are
reported rather than discarded, because they are evidence about the tolerance axis. Where an
experiment was run more than once, a tolerance counts as checked only if it passes in every
run, since separate runs train separate fields. The toy experiments of claims 2 and 3 were
evaluated at a single tolerance, $10^{-6}$, so for them the first rule reduces to checking that
tolerance.

The threshold was fixed before the confirmatory runs, at the same one per cent we require of
agreement between adjoint and direct gradients. The outcomes do not depend on its exact value.
Recomputed at $3\times10^{-3}$ and at $3\times10^{-2}$, every claim keeps its verdict. What
moves is which rung of the ladder a number is quoted at, and how many measurements are excluded.
The extension excludes 25, 6 and 2 of its 300 runs at the three thresholds, and the augmented
model leads the Neural ODE in every cell of the sweep in all three cases.

The consequences are visible in the results below. At $10^{-3}$, the tolerance the original
papers specify, the one-dimensional experiment of claim 1 fails the check for 11 of its 15
trained models, and the worst reconstructs its input with a relative error of 25. In the image experiments the
check fails at the tolerance used for training, on every model and every seed.

\subsection{Deviations from the original descriptions}
\label{sec:deviations}

We reimplemented from the text, and the text does not fix every choice. Where a choice was
unstated we made one and recorded it. Where we departed from a stated choice we list it in
Table~\ref{tab:deviations}, which covers the departures that could move a number. A fuller
register accompanies the code.

One choice is not a departure but needs stating first. We train at a solver tolerance of
$10^{-3}$, the tolerance both original papers specify. The reconstruction check shows that this
tolerance does not integrate the trained fields, so we quote every result at a tolerance that
passes the check rather than at $10^{-3}$. This is where our reporting differs from the
original protocol.

Two departures matter most. We run five seeds against the original twenty repeats on
the toy task, so our intervals are wider and small differences between arms should not be
over-read. For the supporting claims we integrate with an explicit Runge--Kutta method where
Chen et al.\ used implicit Adams. That choice is deliberate, because it is the method a
current user of \texttt{torchdiffeq} will reach for by default, but it is not neutral, and the diagnosis of claim 6 examines it directly.

\begin{table}[h]
  \centering
  \small
  \caption{Departures from the papers as described that could affect a reported number.}
  \label{tab:deviations}
  \begin{tabularx}{\linewidth}{@{}lXX@{}}
    \toprule
    & \textbf{As described} & \textbf{Here, and why} \\
    \midrule
    Reporting tolerance & $10^{-3}$, as for training & the loosest tolerance that passes the
      reconstruction check. Training stays at $10^{-3}$ \\
    Repeats & 20 (toy task) & 5 seeds, and 10 for the skewed held-out metrics. Wider
      intervals \\
    Solver & implicit Adams (Chen et al.) & \texttt{dopri5}, the current default choice.
      Bears on claim 6 \\
    Toy objective & regression to $\pm1$ under MSE & binary classification under
      cross-entropy. The mechanism is unchanged, absolute losses are not comparable \\
    Toy learning rate & $10^{-3}$ & $3\times10^{-3}$ \\
    Parameter counts & 84{,}395 / 84{,}816 (MNIST) & 85{,}316 / 85{,}462, from a flattening
      readout. Matched between arms to within 2\% \\
    Training budget & 50 epochs (toy) & swept over $\{25,\dots,1000\}$, so the original
      budget is one point on a curve \\
    \bottomrule
  \end{tabularx}
\end{table}

\subsection{Experimental setup and code}

Most experiments use five seeds. The held-out metrics of the claim-3 extension are skewed
across seeds, so those use ten. Two use three: claim 8, and the convolutional part of the
claim-6 diagnosis. The peak memory measured for claim 8 does not depend on the seed, so its
three seeds are identical replicates rather than independent samples. We report the mean and standard deviation for image accuracy,
following the original, and the median with interquartile range where the across-seed
distribution is skewed. The statistic used is stated with each result.

Before each run we recorded the condition under which the claim would count as refuted. The
conditions were written down in advance, so that the observation decides the outcome. They are
quoted with the results below, including two cases in which our own expectation was wrong.
Every result row also records the hardware that produced it, which the claim-5 results show was necessary.

The code, the per-seed result files behind every number in this article, and the scripts that
produce every figure are available in the repository cited at the end. One command rebuilds
all figures from the committed results without a GPU. A continuous-integration workflow runs
that path on every commit.

\subsection{Computational requirements}

Experiments ran on one NVIDIA RTX 3090 (24 GB) and on A100 80 GB MIG slices of a SLURM
cluster. The two-dimensional tasks are faster on CPU and ran on an AMD Ryzen 7 7700X. The
matched-parameter MNIST experiment takes about 1.9 hours and CIFAR-10 about 4.3 hours for five
seeds on one A100 slice. The systematic extension of claim 3 takes about 1.6 hours across
fourteen CPU workers. The full set of experiments is on the order of twenty GPU-hours. None of
the claims tested here requires a large compute budget.

Seeding fixes initialisation, data generation and shuffling, so a re-run reproduces our
numbers within the reported spread rather than bit for bit. Adaptive solvers and GPU kernel
selection are not bit-reproducible across machines. Figures rebuilt from the committed result files carry exactly the same
numbers, in the pinned container as on the host. The rendered images are not byte-identical
across environments. Text layout and legend rendering differ between builds of the plotting
stack, while the values drawn do not. One toy experiment re-run two months later, on a different
thread count, returned its held-out accuracies unchanged.

\section{Results}
\label{sec:results}

Each subsection states the claim, the condition for refutation recorded before the run, the
observation at a checked tolerance, and the outcome.

\subsection{Results reproducing Dupont et al.}

\subsubsection{Claim 1: the one-dimensional crossing flow}

Proposition 1 predicts more than failure. A map that preserves order and must send both
classes to opposite constants can do no better than collapsing them onto one value. For
targets $\pm 1$ that gives a mean squared error of exactly 1. This value separates a model
that is structurally obstructed from one that is merely undertrained.

We train with the settings the original gives for $d=1$ and evaluate on a ladder of
tolerances from $10^{-3}$ to $10^{-9}$. Four of the five Neural ODE seeds and every augmented
seed pass the reconstruction check at $10^{-7}$, and we quote the comparison there. The Neural
ODE reaches a median MSE of \textbf{1.0001}, with a range of 1.0001 to 1.0002. Both augmented
models, with $p=1$ and $p=5$, reach \textbf{0.0000}. Forward NFE is 305 against 206 and 170.
The Neural ODE sits at the predicted floor rather than below it, so the claim reproduces in its
strong form.

The fifth Neural ODE seed fails the check at every tolerance we tried. Its reconstruction error
falls steadily as the tolerance tightens, from 25 at $10^{-3}$ to 0.014 at $10^{-9}$, just above
the threshold of 0.01. Its field is stiffer than the others and needs a tighter tolerance than
our ladder reaches. Its MSE is 1.0000 at every tolerance. We report it and leave it out of the
comparison, as we recorded before running the extended ladder.

This experiment also shows the effect of tolerance most clearly. At $10^{-3}$, the tolerance the
original specifies, the reconstruction check fails for 11 of the 15 trained models, and the
worst reconstructs its input with a relative error of 25. The MSE values at that tolerance are
close to those above, so the conclusion is robust, but at that tolerance it was not
established. The augmented models are also the easier ones to integrate. They pass the check on every
seed from $10^{-5}$, while the Neural ODE needs $10^{-7}$ or tighter, because its flow must
steepen to approximate a map it cannot represent.

\resfig{../figures/crossing/crossing_flow.png}
  {Claim 1. Left: median error against solver tolerance. The Neural ODE sits at the error
   floor of 1 predicted by Proposition 1, and the augmented models far below it. Right:
   forward-backward reconstruction error, with the threshold marked. At $10^{-3}$, the
   tolerance the original specifies, most models fail the check.}
  {fig:crossing}

\subsubsection{Claim 2: cost growth with training budget}

We swept the training budget from 25 to 1000 epochs and measured at a checked tolerance of
$10^{-6}$. On the concentric spheres, the median forward NFE of the Neural ODE grows by a factor of
\textbf{1.37} between 25 and 500 epochs, and that of the augmented model by \textbf{1.06}, both
over the four seeds that pass the reconstruction check at both ends. That range of budgets is
the one the condition we recorded in advance names. At the original
budget of 50 epochs, the Neural ODE needs \textbf{218} function evaluations against
\textbf{170} for the augmented model. Both models classify the data almost perfectly at that
budget, with dense accuracies of 0.996 and 1.000. The two models are therefore separated by
the cost of the flow rather than by accuracy. The pattern repeats on the thinner circles geometry, with growth factors of 1.37 and 1.04. The
first equals the spheres figure by coincidence. One of those fifty runs stopped at 271 of its 500 epochs under a wall-clock cap. We include
it, and without it the factor is 1.31.

Three of sixty measurements fail the reconstruction check: one seed at a budget of 200, and
another at 500 and 1000. Their flows become too stiff to integrate at $10^{-6}$, which makes
them the most extreme instances of the effect under study. Over all five seeds, including those three, the factors are 1.74 and 1.04. We report that figure and do not let it carry the comparison, because a measurement that
fails the check is evidence about the tolerance axis rather than about cost.

\resfig{../figures/budget/nfe_vs_epoch.png}
  {Claim 2. Forward NFE against training budget on concentric spheres, at a checked tolerance.
   The cost of the Neural ODE grows during training. The augmented model stays flat and
   lower.}
  {fig:budget}

\subsubsection{Claim 3: generalisation across an unobserved region}

We removed the angular wedge $[0, \pi/5]$ from the training set and tested on that wedge
alone. Over five seeds, all passing the reconstruction check, the Neural ODE reaches a median
accuracy of \textbf{0.619} on the held-out wedge, with a loss of \textbf{6.089}. The augmented
model reaches \textbf{1.000} and \textbf{0.000}. Both models fit the training data perfectly
and both score above 0.96 on an ordinary validation set. The gap appears only on the region
absent from training.

The across-seed spread is large. The held-out accuracy of the Neural ODE has an interquartile
range of 0.613 to 0.921. Most seeds fall to near chance on the wedge, and one seed crosses it
almost perfectly. A mean would misrepresent this distribution, so we report medians, and the
extension below uses ten seeds.

\paragraph{Extension: dependence on the size of the missing region.}
The original demonstrates this claim with one wedge on one geometry. We swept the augmentation
dimension $p \in \{0,1,2,3,5\}$, the wedge width $\{\pi/8, \pi/5, \pi/3\}$ and two geometries
over ten seeds, giving 300 runs. Each run carries the same reconstruction check, and six fail
it and are excluded. The five seeds of the headline run are among these ten, with identical
values. The cell that repeats it reads 0.744 rather than 0.619, because two of the five added
seeds are ones on which the Neural ODE generalises perfectly. The augmented model generalises better at every width, on both
geometries, and for every $p \ge 1$. No condition stated in advance was violated.

The effect does not grow without limit. On the spheres, the advantage in held-out accuracy
increases from $+0.07$ at $\pi/8$ to $+0.26$ at $\pi/5$, then flattens between $+0.27$ and
$+0.30$ at $\pi/3$. The last step lies inside the across-seed spread. Measured by loss rather
than accuracy, the advantage shrinks at the widest wedge, because the augmented model also
degrades there. A large enough hole reduces the performance of both models, and augmentation
does not make an unobserved region free. On the second geometry the gap is smaller and is not
monotone in width. Two further observations follow from the sweep. On the observed region the
median run scores at least 0.998 in every cell, the Neural ODE cells included. Individual runs
dip lower, 27 of 294 below 0.998 and the lowest at 0.818, but these dips are small beside the
held-out failure, which is specific to extrapolation into the gap rather than a general fitting
problem. More augmentation is also
not uniformly better, since $p=2$ generalises best at the widest wedge on both geometries.

\resfig{../figures/slice_grid/slice_grid.png}
  {Extension to claim 3. Held-out accuracy and loss against the width of the removed wedge, by
   augmentation dimension, on two geometries. Median and interquartile range over ten seeds.
   The advantage grows and then flattens, and the augmented model also degrades at the widest
   wedge.}
  {fig:slicegrid}

\subsubsection{Claim 4: matched-parameter image benchmarks}

Table~\ref{tab:table1} gives our matched-parameter results next to the original values. Each
is quoted at the loosest tolerance at which both arms pass the reconstruction check on every
seed, in every run of the experiment we committed. The ordering asserted by the claim reproduces on both datasets. The augmented model is
more accurate by 3.9 percentage points on MNIST and by 6.3 on CIFAR-10. It is also cheaper,
requiring 1.91 times fewer function evaluations on MNIST and 1.26 to 1.32 times fewer on
CIFAR-10, at the tolerances that pass the check in every run. The ratio is less stable across
runs than accuracy: an independent retrain of the MNIST models gives 2.10 at $10^{-7}$.

Two of the four absolute values match the original and two fall below it. The CIFAR-10 Neural
ODE matches within noise, as does the MNIST augmented model. The MNIST Neural ODE falls about
2.2 points short, consistently across two independent runs. The CIFAR-10 augmented model falls
0.6 points short in the run reported here and 1.3 points short in an earlier run on different
hardware. That spread is larger than the shortfall, so we report a small undershoot rather
than a match or a discrepancy. We did not tune towards the published values.

The MNIST Neural ODE shortfall needs its own explanation, because the augmented model,
trained under identical settings, matches. Seed noise does not explain it: the gap is about
five standard deviations and repeats across two runs on different hardware. Our hypothesis is
that eight epochs under-train the Neural ODE but not the augmented model. At the final epoch
the Neural ODE reaches a training accuracy of only 94.3\%, against 98.6\% for the augmented
model, so it is underfitting rather than overfitting. With parameter counts matched, that points to
optimisation rather than capacity. Its test accuracy also rises by 0.5 points over
the last epoch in both runs, while the augmented model's is flat from the seventh epoch. The
original states neither the epoch count nor the learning-rate schedule for these experiments,
so a longer budget or a decaying schedule could close the gap for the Neural ODE alone.

We tested the budget with three seeds trained for 20 epochs, otherwise unchanged. Before the
run we required the epoch-20 accuracy, at a tolerance that passes the reconstruction check on
all three seeds, to reach 95.9. That test could not be applied: by epoch 20 two of the three
fields had stiffened beyond $10^{-7}$, the tightest tolerance on our ladder. The trajectories
support the hypothesis without confirming it. At the training tolerance the mean test accuracy
passes 95.9 at epoch 15 and peaks at 96.1 at epoch 17, with all three seeds between 96.0 and 96.3, inside the original's interval. Training accuracy
rises to 96.5. One seed then
diverged and ended at 92.3, while the other two ended at 96.1 and 96.2. Without a learning-rate
schedule the longer run is unstable. We therefore attribute the undershoot to the training
budget with moderate confidence, and keep the 8-epoch value in Table~\ref{tab:table1}. The conditions recorded before the run concerned the ordering, which holds.

\begin{table}[h]
  \centering
  \small
  \caption{Matched-parameter image results, mean $\pm$ s.d. over five seeds, quoted at the
           loosest tolerance at which both arms pass the reconstruction check on every seed,
           in every committed run.}
  \label{tab:table1}
  \begin{tabular}{@{}llccc@{}}
    \toprule
    Dataset & Model & This replication & Original & Tolerance \\
    \midrule
    MNIST    & Neural ODE & $94.16 \pm 0.44$ & $96.4 \pm 0.5$ & $10^{-7}$ \\
    MNIST    & ANODE ($p=5$)  & $98.05 \pm 0.20$ & $98.2 \pm 0.1$ & $10^{-7}$ \\
    CIFAR-10 & Neural ODE & $53.70 \pm 0.83$ & $53.7 \pm 0.2$ & $10^{-6}$ \\
    CIFAR-10 & ANODE ($p=10$) & $60.04 \pm 0.94$ & $60.6 \pm 0.4$ & $10^{-6}$ \\
    \bottomrule
  \end{tabular}
\end{table}

\paragraph{The tolerance at which accuracy is measured.}
These models are trained at $10^{-3}$, and at that tolerance the reconstruction check fails
for every model and every seed on both datasets. Our first version of this experiment measured
accuracy there, so the reported values came from a regime in which the solver was not tracking
the flow. We repeated both experiments with the training tolerance included in the evaluation
ladder and the test set re-evaluated at every tolerance. Accuracy is almost insensitive to the
change. It moves by at most 0.11 points on average, and by 0.40 points for the worst single
seed, between the training tolerance and a checked one. Cost behaves differently, since the
forward NFE at $10^{-7}$ is several times its value at $10^{-3}$. We therefore quote both
accuracy and cost at checked tolerances.

\subsection{Results reproducing Chen et al.}

\subsubsection{Claim 5: solver dynamics}
\label{sec:claim5}

We trained a convolutional ODE-Net, froze it, and swept the solver at evaluation. For each
configuration we measured the endpoint error against a high-accuracy reference, an
eighth-order method at $10^{-8}$, together with the forward wall-clock time and the NFE.

Both parts of the claim hold. The error falls monotonically as the tolerance tightens, for
every adaptive solver tested. Forward time is strongly rank-correlated with NFE, with
Spearman's $\rho$ between 0.983 and 0.992 per solver.

The cost part of the claim was recorded in advance as monotone, and by that condition it is
violated. The median NFE of \texttt{dopri5} falls from 26 to 20 between tolerances $10^{-1}$
and $10^{-2}$. We report the violation as measured rather than restating the condition. Both
of those measurements fail the reconstruction check on every seed, so the step sizes chosen
there do not track the flow. Among measurements that pass the check, the cost is monotone.
Cost therefore rises as the tolerance tightens wherever the solver is integrating, and behaves
irregularly where it is not.

Two further observations follow from this experiment. First, the cost of the reference grows by a factor of between 4.7 and 7.7 per decade of tolerance, where an eighth-order method on a non-stiff
field would cost about 1.3 times per decade. This is a direct measurement of the stiffness of
a trained field, and it explains why high-accuracy references on these models are expensive.
Second, a nominal tolerance is not comparable across solvers. At $10^{-5}$, \texttt{dopri5}
uses 116 function evaluations and sits at the edge of the reconstruction check, while
\texttt{bosh3} uses 227 and \texttt{dopri8} uses 1055 for reconstruction errors an order of
magnitude smaller. The tolerance is a request rather than a guarantee, and the reconstruction
check is what makes measurements comparable.

The reference itself is not exact, and we measured how inexact. On one field retrained here
with the same configuration, its endpoint
at $10^{-8}$ differs from its endpoint at $10^{-9}$ by a relative $6\times10^{-6}$. That is
below a tenth of 95\% of the errors on this axis, and between 13\% and 22\% of the four
tightest. The monotone fall of the error, which spans four orders of magnitude, is therefore
unaffected, while the smallest errors we report carry a reference uncertainty of about a fifth
of their value. Tightening the reference further is a matter of cost rather than difficulty:
209{,}783 function evaluations at $10^{-8}$ against 1{,}033{,}060 at $10^{-9}$, a factor of 4.9
per decade, in line with the growth reported above.

This experiment also required the hardware recorded with each measurement. The five seeds were
scheduled across two different GPU partition sizes, so the pooled wall-clock times mixed
machines. Within each hardware group the correlation between time and NFE is between 0.991 and
0.997, which is tighter than the pooled value. The confound weakened the result rather than
producing it. This check is part of the analysis script.

\resfig{../figures/c1/solver_dynamics.png}
  {Claim 5. Left: endpoint error against tolerance for the adaptive solvers. Right: forward
   time against function evaluations, for adaptive and fixed-step solvers. Cost is linear in
   NFE.}
  {fig:solver}

\subsubsection{Claim 6: backward versus forward NFE}
\label{sec:claim6}

This is the one claim we do not reproduce. Rather than report the discrepancy alone, we
measured its cause.

On a trained convolutional ODE-Net, at a tolerance that passes the reconstruction check, the
backward pass costs \textbf{123 times} the forward pass in function evaluations with
\texttt{dopri5}, and 203 times with \texttt{bosh3}. The original reports about one half, but its own logs, discussed below, give about 0.8. On
the two-dimensional fields the ratios are smaller but still far above one, between 25 and 96
for \texttt{dopri5}. The gradients are correct in all cases. Every value quoted here comes
from a configuration whose parameter gradients agree with direct backpropagation at high
accuracy to within 1\%, so the cost is real and not the result of a failed solve.

We tested three explanations, each with the effect size that would count recorded in advance.
Of 456 configurations, 18 did not
complete: twelve failed at once, for memory, on the convolutional field, and six reached a
ten-minute cap on the toy field. 316 passed both the reconstruction check and the gradient
check. Only those carry the conclusions.

\textbf{Solver family.} The original used implicit Adams, a stiff-capable method, while we
used an explicit Runge--Kutta pair. The adjoint integrates an augmented reverse system that is
stiffer than the forward one, so this is the natural explanation. It is wrong where we could test it. On the toy field, \texttt{scipy:LSODA} gives a ratio of
2159 and \texttt{scipy:BDF} gives 28, against 25 for \texttt{dopri5}. On the convolutional field
neither could be run at all: \texttt{scipy:LSODA} fails immediately, attempting a dense Jacobian
of the 401{,}408-dimensional state, about 1.2 TiB, and \texttt{scipy:BDF} was not attempted for
the same reason. The stiff-solver explanation is therefore refuted on the toy field and untested
on the convolutional one.

We later ran the original's own solver. \texttt{torchdiffeq} cannot drive VODE, so we wrote its adjoint directly, following the configuration of the
code the authors shared. The augmented reverse system is integrated backwards by VODE's implicit
Adams method at the forward tolerance, with the field evaluated in \texttt{torch}. Its gradients
are checked against direct backpropagation. On three convolutional fields retrained here with the
configuration of the committed run it gives a median ratio of \textbf{108.9}, against 128.4 for
\texttt{dopri5} on those same fields. On the two-dimensional field it gives \textbf{9.0} against
24.6. It is the more expensive solver in both directions. It moves the ratio by a factor of 1.2
on the convolutional field and 2.7 on the toy one, where the gap to the original is a factor of
about 130. The solver family is not the explanation either.
At the looser tolerances of the original's figure its gradients fail our one-percent check on
both fields, so that part of the range cannot be compared here.

\textbf{Adjoint tolerance.} This is the dominant factor. The reverse system need not be solved
at the tolerance of the forward pass, but \texttt{torchdiffeq} couples the two by default, and
so did we. Solving the adjoint at a tolerance 100 times looser reduces the ratio by a factor
between 5 and 76, across 15 of 17 solver and tolerance combinations. On the convolutional
field it takes the ratio from 123 to \textbf{4.1}, with gradients still correct to within 1\%.
Within our configuration the adjoint tolerance is therefore the dominant factor. It does not
explain the gap to the original, which also solved the reverse system at the forward tolerance.

\textbf{Field stiffness.} We expected the ratio to grow as the field trains, and recorded in
advance that it would do so for every solver. It does not. In 5 of 16 combinations the
untrained field has the higher ratio, so this explanation is refuted as stated. It also
corrects an earlier claim of our own. The reason appears to be that the ratio combines two
quantities. On a field that is easy to integrate the forward count is small, so a moderate
backward cost still produces a large ratio. A ratio is a poor summary statistic for this
quantity, and the two counts are better reported separately.

\textbf{The original measurement.} After our experiments for this claim had finished, the
original authors shared the code and logs behind Figure~3 (Section~\ref{sec:communication}). The reverse system was solved with the implicit Adams method of SciPy's VODE, at the same
tolerance as the forward pass, as in our default configuration. The backward count is the number
of evaluations of the augmented system, as in ours. The forward counts plotted in Figure~3c,
however, are not forward counts alone. The evaluation counter is reset after the forward count is
recorded, and the backward pass then increments it. Each recorded forward count therefore
contains the previous batch's backward count as well, plus one evaluation that the adjoint makes
outside the reverse solve. Table~\ref{tab:chen3c} sets out the counts. In every row the median
recorded forward count is the corrected forward count plus the backward count plus one, to
within one evaluation. At a tolerance of 1, where the counts are constant, every backward pass
takes 6 evaluations and every recorded forward count after the first is exactly $13 = 6 + 6 + 1$.
Where the counts vary most, at the three tightest tolerances, each recorded forward count
correlates with the previous batch's backward count rather than with its own. A backward count
somewhat below the forward count, plotted against their sum, appears as about one half
(Figure~\ref{fig:chen3c}). Corrected, the original's ratio is \textbf{0.78 to 0.86} between $10^{-5}$ and $10^{-1}$,
and 1.00 at a tolerance of 1, where the counts are smallest.

\begin{table}[h]
  \centering
  \small
  \caption{Claim 6. Evaluation counts in the original logs behind Figure~3c, medians over the
           101 batches at each tolerance. $B$ is the backward count and $L$ the recorded forward
           count, which the original plots. The corrected forward count is
           $F_k = L_k - B_{k-1} - 1$, and $L_0/2$ for the first batch, whose forward pass runs
           twice; $L$ and $B/L$ exclude that batch. The last two columns correlate $L_k$ with the
           previous and with the same batch's backward count, and are undefined where the
           backward count is constant. From \texttt{scripts/chen\_fig3c\_audit.py}.}
  \label{tab:chen3c}
  \begin{tabular}{@{}lrrrcccc@{}}
    \toprule
    Tolerance & $B$ & $L$ & $F$ & $B/L$, plotted & $B/F$, corrected & $r(L_k, B_{k-1})$ & $r(L_k, B_k)$ \\
    \midrule
    $10^{-5}$ & 50 & 116 & 64 & 0.44 & 0.79 & 0.77 & 0.25 \\
    $10^{-4}$ & 26 &  60 & 32 & 0.44 & 0.81 & 0.61 & 0.26 \\
    $10^{-3}$ & 15 &  36 & 20 & 0.43 & 0.78 & 0.56 & 0.29 \\
    $10^{-2}$ &  9 &  21 & 11 & 0.43 & 0.82 & 0.10 & 0.05 \\
    $10^{-1}$ &  6 &  14 &  7 & 0.43 & 0.86 & --   & --   \\
    $1$       &  6 &  13 &  6 & 0.46 & 1.00 & --   & --   \\
    \bottomrule
  \end{tabular}
\end{table}

The original authors have reviewed this analysis and agree with it
(D. Duvenaud, personal communication, 11 September 2026). The original took the ratio of one half as
evidence that the adjoint is cheaper than backpropagating through the solver. At 0.8 that argument still holds for the original's configuration, by a smaller margin. In ours
the reverse pass costs 13 to 123 times the forward pass.

\resfig{../figures/chen_fig3c/fig3c_rebuilt.png}
  {Claim 6. Panel (c) of the original Figure 3, rebuilt from the per-batch counts in the
   original logs, which the authors shared. Left: backward count against the logged forward
   count, as in the original figure. Right: against the forward count with the previous
   batch's backward count removed. Colours follow the original, from a tolerance of 1 (red)
   to $10^{-5}$ (purple).}
  {fig:chen3c}

\textbf{Error norm.} One further difference between the two configurations was a candidate.
VODE measures local error with a single weighted root-mean-square norm over the whole augmented
state, which includes the adjoint of every parameter. \texttt{torchdiffeq} takes the maximum of
a separate norm for the state, its adjoint and each parameter tensor, which is stricter at the
same tolerance. We tested this on the convolutional field, with the adjoint at the forward
tolerance of $10^{-5}$, recording in advance that a reduction below two times would refute it.
A VODE-style norm reduces the ratio from 128 to 79, by 1.6 times, and the seminorm of
\citet{kidger2021hey}, which leaves the parameter adjoints out of the norm, reduces it to 98.
The explanation is refuted. Two of the three retrained fields pass both checks. The third
narrowly fails the reconstruction check at this tolerance and is excluded, although its unchecked
ratios fall by a similar factor. The gap between our ratio and the original's
corrected value of about 0.8 therefore remains unexplained: neither the adjoint tolerance,
nor the solver, including the original's own, nor the error norm accounts for it.

The original's corrected value is reachable in a checked configuration. The lowest ratio
we observe is 0.78. Reaching it requires solving the adjoint at a looser tolerance than the
forward pass, and it occurs on the toy field with a low-order solver. The library default is to solve the adjoint at the forward tolerance, which is also the natural
reading of the original setup. Under that setting the lowest ratio we observe anywhere is 5.96,
and on our convolutional field it is 119. We therefore report claim 6 as not reproduced, with the qualification
that the quantity is a property of the configuration of the adjoint rather than of the adjoint
method. 

One limitation bounds this conclusion, and it is now the field itself. Three named differences between the two configurations have been ruled out: the adjoint
tolerance, the solver and the error norm. Both configurations couple the adjoint tolerance to the
forward tolerance, and the solver we ran is the original's own. The
original's network is smaller, group-normalised and trained for 120 epochs at a tolerance of
$10^{-2}$, while ours trains for five epochs at $10^{-3}$. A reverse pass costing about a
hundred times the forward therefore looks like a property of the trained field rather than of
the adjoint method. We did not vary the architecture and the training budget to test that, and
we do not claim it.

\resfig{../figures/c2_diagnosis/c2_diagnosis.png}
  {Claim 6. Backward over forward NFE against forward tolerance, on the toy field (left) and
   the convolutional field (right). Solid lines solve the adjoint at the forward tolerance,
   dashed lines 100 times looser. The dotted line is the value reported in the original. Only
   configurations passing both the reconstruction and gradient checks are shown.}
  {fig:c2diag}

\subsubsection{Claim 7: NFE growth during training}

On the convolutional ODE-Net, at $10^{-7}$, a tolerance that passes the reconstruction check on
every seed and epoch, the forward NFE grows from \textbf{384 to 738} over six epochs. This
holds in all five seeds, so the claim reproduces.

The behaviour of the tolerance is more informative than the NFE itself. Over the same six
epochs, the fraction of seeds passing the check at $10^{-5}$ falls from 1.0 in the first three epochs to between 0.4 and 0.6 in the last three. At
$10^{-6}$ and $10^{-7}$ every seed continues to pass.
The loosest tolerance that passes on every seed tightens from $10^{-5}$ to $10^{-6}$ at the
fourth epoch. The field does not only become more expensive to
integrate. It becomes harder to integrate at a fixed tolerance. This is the mechanism behind
Section~\ref{sec:recon}, measured directly, and it is why a tolerance chosen at the start of
training cannot be assumed to remain adequate. A companion ten-epoch run at a single tolerance
of $10^{-5}$ shows the consequence. Its NFE curve rises smoothly and looks reasonable, but by
the final epoch none of its five seeds passes the check, so the second half of that curve does
not describe an integrated flow.

\resfig{../figures/mnist_stiffening/stiffening.png}
  {Claim 7. Forward NFE over training at a checked tolerance, and the fraction of seeds passing
   the check at each tolerance. The field becomes both more expensive and harder to integrate.}
  {fig:stiffening}

\subsubsection{Claim 8: constant-memory adjoint}

The memory claim concerns the effective depth of the ODE-Net itself. Varying the depth of a
discrete baseline and reporting the ODE-Net as a single point measures a different quantity,
which is what an earlier version of this project did. Testing the claim requires holding the
model fixed and increasing its own number of function evaluations.

We did this by tightening the solver tolerance. Peak memory under the adjoint is flat at
\textbf{$+0.0007$ MB per function evaluation}, staying at 442 MB as the NFE rises from 14 to 92. Direct backpropagation through the solver rises at \textbf{$+31.8$ MB per function
evaluation}, from 689 MB to 3185 MB over the same range. The claim reproduces. The gradients of
the two methods agree throughout to within half a percent, so the constant memory does not come
at the cost of a worse gradient. Peak memory does not depend on the seed, and the three seeds we ran give identical values.
They show that the measurement repeats, not that it holds across independent samples.

\resfig{../figures/c4/memory_vs_nfe.png}
  {Claim 8. Peak memory against the forward NFE of the model, increased by tightening the
   solver tolerance. Memory is flat under the adjoint and linear under direct
   backpropagation.}
  {fig:memory}

\section{Discussion}
\label{sec:discussion}

Table~\ref{tab:summary} collects the outcomes. Five of the eight claims reproduce, two reproduce
partially, and one is not reproduced. Both partial cases keep the structure of the original
claim and miss on a number. On the image benchmarks the ordering and the cost advantage
reproduce on both datasets, while two of the four absolute accuracies fall short. On the solver
dynamics both substantive relations hold, while a condition we recorded as monotone is violated
at tolerances where the solver is not integrating. Neither case is a failure of the underlying
claim, and neither is a clean success.

The structural claims of the original paper are robust. They follow from the topology of a
flow, they appeared immediately, and they survived every tolerance, geometry and seed we tried.
In the one-dimensional case the model fails by exactly the margin the theory predicts. The
claims about cost are the fragile ones. The fragility arises from the interaction between the learned field
and the solver.

\begin{table}[h]
  \centering
  \small
  \caption{Outcome per claim and the evidence it rests on. Every quantity is quoted at a
           tolerance at which the reconstruction check passes, except the all-seed growth
           factor of claim 2, which is reported beside the checked one and includes seeds that
           fail at the largest budgets.}
  \label{tab:summary}
  \begin{tabularx}{\linewidth}{@{}clX@{}}
    \toprule
    \# & Outcome & Evidence \\
    \midrule
    1 & reproduced & Neural ODE at the predicted error floor, 1.0001 against 0.0000, on four of
        five seeds. The fifth is too stiff to check and also sits at the floor \\
    2 & reproduced & NFE growth of $\times1.37$ against $\times1.06$ on the seeds that
        pass the check at both ends, and $\times1.74$ against $\times1.04$ over all five \\
    3 & reproduced & held-out accuracy 0.619 against 1.000, extended over widths and
        geometries \\
    4 & partial & ordering reproduced on both datasets. Two of four accuracies fall short by
        0.6 and 2.2 points \\
    5 & partial & error and time-NFE relations hold. The monotone cost condition is violated
        at tolerances where the solver is not integrating \\
    6 & not reproduced & 123 against the original's corrected value of about 0.8
        (reported as 0.5) \\
    7 & reproduced & NFE 384 to 738 over six epochs, and the checked tolerance tightens during
        training \\
    8 & reproduced & $+0.0007$ against $+31.8$ MB per function evaluation \\
    \bottomrule
  \end{tabularx}
\end{table}

\subsection{What we withdrew}
\label{sec:withdrawn}

Several conclusions were recorded during this work and later withdrawn. We list them because
the record of what failed a check is part of the evidence for what passed one.
Table~\ref{tab:withdrawn} gives each in the order it happened.

\begin{table}[h]
  \centering
  \small
  \caption{Conclusions withdrawn during this work, in order, and what replaced them.}
  \label{tab:withdrawn}
  \begin{tabularx}{\linewidth}{@{}XXX@{}}
    \toprule
    First recorded & Why it was withdrawn & What replaced it \\
    \midrule
    The Neural ODE does not fail on the toy task, and its NFE stays flat at about 44 &
      Measured at $10^{-3}$, where the flow is not integrated. Reconstruction error was about
      0.9 on data of radius about 1, and a linear probe on the terminal state fell to chance &
      Re-measured at $10^{-6}$: NODE cost grows with training, as the original reports (claim 2) \\
    Backward NFE is about equal to forward NFE (claim 6) &
      The value came from a loose tolerance. Tightening it moved the ratio from 25 to 115 &
      No single ratio exists. Claim 6 is reported as not reproduced, with a diagnosis \\
    On MNIST the augmented model is 1.7 times cheaper in NFE (claim 4) &
      Measured at $10^{-5}$, where three of five Neural ODE seeds fail the check &
      2.10 at $10^{-7}$ in a retrain, 1.91 in the reported run \\
    The Neural ODE cannot represent the crossing map (claim 1) &
      Measured at $10^{-3}$, where most models fail the check. The conclusion was not
      established there, although the numbers did not move &
      Re-measured on a ladder to $10^{-9}$ at the original settings. It holds at $10^{-7}$ \\
    Adjoint memory is constant in depth (claim 8) &
      We varied the depth of a discrete baseline, which tests a different quantity &
      Hold the ODE-Net fixed and raise its own NFE \\
    \bottomrule
  \end{tabularx}
\end{table}

The corrections did not all run in one direction. Withdrawing the flat NFE made the original
result reproduce, and re-measuring the MNIST cost ratio enlarged the advantage of the augmented
model. Withdrawing the equal-cost ratio took away what had looked like a clean headline, and
left a claim we do not reproduce. We did not choose among them by outcome. Each was decided by a
check, and four of the five by the same one: whether the solver had tracked the flow.

Smaller errors in our own documentation were corrected in the same way. An earlier deviation
register gave the removed wedge as $[0,\pi/13]$, which appears nowhere in the original, and
reported medians as means. Neither changed a result.

\subsection{What was easy}

The structural claims are robust and appear immediately. The one-dimensional failure and the
missing-region gap were visible in the first run of each experiment, at every tolerance, and
survived every later change of protocol. The appendices of both papers specify the
architectures precisely enough to reimplement without guesswork. \texttt{torchdiffeq} makes the
solver side straightforward, including the adjoint. The toy experiments are cheap. The whole
two-dimensional programme, including a 300-run extension, runs in a few CPU-hours.

\subsection{What was difficult}

Solver tolerance was the main obstacle, and it is difficult because nothing fails. The
integrator returns a state, the counter returns a number, and only an explicit check shows that
the flow was not tracked. We added the reconstruction check after reviewing a set of results
that looked unusually clean. It invalidated several conclusions we had already recorded, listed in
Section~\ref{sec:withdrawn}. A tolerance should be treated
as unverified until checked, whether it is a library default or the value an original paper
used, and it should be rechecked as training proceeds.

Three further difficulties arose. Epoch budgets are not stated for the image
experiments in the original, so our undershoot on two accuracies could not be attributed cleanly. A longer run points to the
budget for the MNIST Neural ODE but could not confirm it, because its fields stiffen beyond our
tolerance ladder.
High-accuracy references are expensive on trained fields, for the stiffness reason quantified
under claim 5, which limits how tightly the error axis of a solver experiment can be pinned.
The metrics of the generalisation experiment are skewed across seeds, since most Neural ODE
seeds fail on the held-out wedge while an occasional one does not. Five seeds are enough to see
that effect but not to characterise it, and a one-seed pilot of the extension pointed in the
wrong direction.

A final difficulty was our own infrastructure rather than the models. The harness for the
CIFAR-10 experiment deleted its results when restarted. The released code writes per-seed files
and resumes from completed ones.

\subsection{Communication with the original authors}
\label{sec:communication}

Claim 6 is the one we report as not reproduced. Before attributing the difference to the original, we measured its cause. The diagnosis tests three explanations against thresholds recorded in advance, and
identifies the coupling between the adjoint and forward tolerances as the dominant one.

Two questions remained that only the original authors could answer. The first is the tolerance
at which the reverse solve of Figure~3c was run, relative to the forward pass. The second is
whether its backward count covers evaluations of the vector field or steps of the augmented
adjoint system. We wrote to
them on 10 September 2026. David Duvenaud and Ricky Chen replied on 11 September and shared the original code
and logs for Figure~3, which are not public. We read them after the diagnosis of claim 6 was complete. The two
experiments that follow it, on the error norm and on the original's solver, were designed in
response to what the code showed. No model or training code here derives from
theirs. One diagnostic script does follow the adjoint in their integrator, so that claim 6 can
be measured in the original's configuration, and \texttt{PROVENANCE.md} records it. The reverse solve ran at the forward tolerance, and the backward count is the number
of evaluations of the augmented system. The logs also showed the counting artefact described under claim 6. We sent
them our analysis, and they agreed with it and plan to update their paper. We are grateful for
their openness in sharing unreleased code, without which this correction would not have been
possible.

\section{Conclusion}

The central argument of Augmented Neural ODEs replicates. A Neural ODE cannot represent the
one-dimensional crossing map, and it fails at the error the proposition predicts rather than
merely performing badly. Its integration cost grows during training, while the cost of an
augmented model does not. It generalises poorly into a region absent from training, and
augmentation repairs this at every width and on both geometries we tested. At matched parameter
counts the augmented model is more accurate and cheaper on both image benchmarks, although two
of the four absolute accuracies fall below the published values. Of the four supporting claims
taken from Chen et al., two reproduce, one reproduces partially, and one is not reproduced.

The results that proved fragile are those about cost. Cost is a property of the solver as much
as of the model. An adaptive solver at too loose a tolerance returns an answer and an evaluation
count, and reports nothing about whether it tracked the flow. The field of a trained Neural ODE
also stiffens during training, so a tolerance fixed at the start stops being adequate. Most of
the conclusions we withdrew during this work follow from this, as does the one claim we do not
reproduce. We suggest four practices for measuring these models:

\begin{itemize}
  \item attach an invertibility check to every measurement,
  \item treat the solver tolerance as an experimental axis,
  \item report the tolerance that each number came from,
  \item report forward and backward evaluation counts separately rather than as a ratio.
\end{itemize}

None of this requires a large budget. The programme reported here, including a 300-run
extension of one experiment, is on the order of twenty GPU-hours. The claims of both original papers are cheap to check. They require one addition: an explicit
check that the differential equation was solved.

\section*{Code and data availability}

The implementation, the per-seed result files behind every number in this article, and the
scripts that produce every figure are available at
\begin{center}\url{https://github.com/VanniLeonardo/re-augmented-neural-odes}\end{center}
The per-seed files are the committed
artifact and the figures are not. One command rebuilds all figures from those files in about
two minutes without a GPU, and prints with each one the condition recorded in advance for the
claim it supports. A continuous-integration workflow builds the pinned container from scratch
and runs that path on every commit. MNIST and CIFAR-10 are fetched from their original sources
at run time and are not redistributed.

\section*{Author contributions}

All authors contributed equally to this work. Authors are listed in alphabetical order
of surname, and the order does not reflect the size of each contribution.

\section*{Acknowledgements}

This replication grew out of a course project at Bocconi University, and the GPU experiments
ran on the university high-performance computing cluster. We thank the authors of both original
papers for work specified precisely enough to be checked, and David Duvenaud and Ricky Chen for answering our questions about
claim 6 and sharing the code behind it.

Generative AI tools were used during preparation of the manuscript and repository documentation
for language editing, restructuring, and drafting assistance. All source code, experiments,
analyses, scientific claims, and conclusions were produced and verified by the authors, who take
full responsibility for the final text.
